\scanpage{209}
rôle de \uline{classes de chemins}, ils sont déterminés \inserted{(s'il y en a au moins un)} \unsure{mod} automorphismes des premiers, qui sont les pts d'un groupe algébrique. Mais en général il peut y avoir des pts à valeurs dans un même $h$ non isomorphes. Les groupes de Galois opèrent sur \deleted{les points} l'ens des pts à valeurs dans $h$. Si on a un pt $\xi$ à val dans $h$, tous les autres pts s'en déduisent par des fibrés principaux \unsure{hom.} sous le (pro)groupe algébrique $\underline{\operatorname{Aut}}(\xi)$.

\diagram{\begin{gathered}
\left\{\begin{array}{ll}
\text{Pts }l\text{-adiques}&\text{valeurs }\mathbf Z_l\text{ ou }\mathbf Q_l\\[3pt]
\text{Pts transcendants}&\begin{gathered}\text{valeurs }\mathbf R\ \deleted{\text{ou }\mathbf C\ \unsure{\text{top}}}\text{ voire }\mathbf Z\\
\text{[structure supplémentaire]}\\
\text{Il semble que ces pts correspondent}\\
\text{biunivoquement aux pts de}\\
X\text{ à valeurs dans }\mathbf C,\\
\text{mod conjugaison}
\end{gathered}
\end{array}\right.\\[6pt]
\begin{array}{ll}
\text{Pts Hodge-de Rham}&\text{valeurs }h\text{ quelconque}\\
\text{pts de Witt}&\text{valeurs }W_p(k)\ \inserted{\text{$k$ parf de car $p$}}
\end{array}
\end{gathered}}

\mspage{208}
